\clearpage
\nobalance
\appendix

\section{Additional experimental details}

\subsection{Benchmark domains}
\label{app:dataset-details}

The ten graph benchmarks comprise social networks (Facebook, GPlus, Twitter,
Flickr, YouTube, and LiveJournal), a trust network (Epinions), citation networks (Cit-HepPh
and Cit-HepTh), and a co-authorship network (DBLP). Vertices correspond to
users, papers, and authors, respectively; edges encode social or trust links,
citations, and co-authorship relations. They represent the graph after entity
and relation extraction. Directed source relations are symmetrized, and all
experiments process undirected graphs. MultiHop-RAG instead supplies the
document-ingestion and downstream question-answering workloads.

\subsection{Execution and accounting defaults}
\label{app:exp-defaults}

% Main-table configuration: source deferral enabled, epsilon=0.05.
% Component-ablation values retain their archived partition-derived cost.
% The revised epsilon sweep separately reports executed schedule cost.
Structural weights count entities at leaves and children at internal nodes;
the complete full-index rebuild normalizes their sum to $1$. Main-table
maintenance latency excludes batch parsing and partition output; Exp-2b
driver timings include index emission. The main-table configuration enables
source deferral with $\varepsilon=0.05$. The paired deferral comparison
uses the same threshold; Exp-5 retains its archived eager-source configuration.
Main-table structural scores and schedule-plus-balance accrual costs are
distinguished in Section~\ref{sec:exp:setup}; partition changes alone do not
recover a deferred execution schedule.

All reported \DyLo{} runs use the current marginal bucket test of
Section~\ref{sec:opt:pruning}, with $\bar w_t(B)$ rather than the static weight
$w(B)$. Across the 18 \DyLo{} detector--graph configurations in the six jointly
measured columns of Table~\ref{tab:main}, the seed-to-seed latency range reaches
$48\%$ on the smallest graph but stays below $5\%$ on the largest, where
sub-millisecond timing noise matters less.

LiveJournal ($4.8$M vertices, $69.0$M edges) is run once per method over the
first $200$ of its $6{,}900$ batches; from-scratch Leiden alone needs $74$ hours
for these batches, and from-scratch \Louvain{} and SLM were not run. Baseline
structural costs are recomputed from the saved per-batch partitions. Peak
memory is $9.5$\,GB for all three \DyLo{} instances and $4.2$--$54.5$\,GB for
the baselines.

\subsection{Motivating measurements}
\label{app:motivation}

Figure~\ref{fig:span} and Figure~\ref{fig:share} replay HIT-Leiden's per-batch
partitions on the $|\Delta E|=10^3$ streams of Table~\ref{tab:main}, using the
first $40$ batches of each graph (all batches of Facebook and Cit-HepTh). Because the baseline stores only a flat partition, we build its
hierarchy by repeatedly contracting the leaf partition, as GraphRAG does when
it refines a flat partition hierarchically; $w(c)$ follows
Definition~\ref{def:cost}, and $\mathrm{Anc}(c)$ is the full root path.

\emph{Recurring refreshes.} Under simulated eager refresh, every source
community that loses a member in a batch refreshes its root path. A
source-side refresh of node $c$ is recurring if $c$ was already refreshed for
the same reason in an earlier batch; Figure~\ref{fig:span} reports the
recurring share of source-side refresh cost.

\emph{Shared refreshes.} For the moves accepted in a batch,
Figure~\ref{fig:share} divides
$\sum_a w\bigl(\mathrm{Anc}(\mathrm{src}(a))\cup\mathrm{Anc}(\mathrm{dst}(a))\bigr)$
by $w\bigl(\bigcup_a\mathrm{Anc}(\mathrm{src}(a))\cup\mathrm{Anc}(\mathrm{dst}(a))\bigr)$,
which is at least one, and averages over batches; the axis is logarithmic
because the ratio spans three orders of magnitude.

\emph{Locality of flipped decisions.} On ten sampled batches per graph, we
compute each vertex's best-move raw gain on the pre-batch partition, apply the
batch with the partition fixed, update $m$ and the community volumes, and
recompute the gain. On every graph, no vertex whose gain turns from
non-positive to positive lies more than two hops from an inserted edge. Gains
of untouched vertices change only through $m$ and the community volumes, by
amounts too small to flip their sign in these measurements. A variant that asks
which vertices of the post-batch graph prefer to move under the old partition
finds $0.0$--$0.9\%$ of them beyond two hops.

\subsection{Certified-frontier statistics}
\label{app:frontier}

Table~\ref{tab:frontier} reports per-batch frontier sizes and certificate
events under the configuration of Table~\ref{tab:main} ($\lambda=10^{-1}$),
over the first $40$ batches of each graph (all batches of Facebook and
Cit-HepTh). The frontier counts the vertices re-examined in a batch; the skip
rate is the share of certificate checks that end without a recheck.

\begin{table}[tb]\centering
\caption{Certified frontier and certificate events per batch.}
\label{tab:frontier}
\begin{footnotesize}
\setlength{\tabcolsep}{0.3em}
\begin{tabular*}{\linewidth}{@{\extracolsep{\fill}}lrrrrr@{}}
\toprule
Graph & Frontier (med.) & Range & Expired & Skip (\%) & Certificates \\
\midrule
Facebook  & $3$   & $0$--$9$     & $0$  & $100.0$ & $2{,}662$ \\
Cit-HepTh & $24$  & $0$--$33$    & $2$  & $99.8$  & $53{,}227$ \\
Cit-HepPh & $20$  & $10$--$34$   & $4$  & $99.8$  & $48{,}149$ \\
Epinions  & $134$ & $117$--$154$ & $18$ & $99.1$  & $84{,}885$ \\
DBLP      & $90$  & $72$--$111$  & $1$  & $100.0$ & $180{,}505$ \\
Twitter   & $11$  & $5$--$17$    & $0$  & $100.0$ & $92{,}825$ \\
Flickr    & $103$ & $85$--$138$  & $8$  & $99.6$  & $306{,}354$ \\
YouTube   & $277$ & $244$--$315$ & $12$ & $99.4$  & $494{,}420$ \\
GPlus     & $2$   & $0$--$6$     & $2$  & $99.9$  & $347{,}580$ \\
\bottomrule
\end{tabular*}
\par\smallskip\raggedright
Expired: certificates expired per batch; Certificates: live certificates after
the last batch.
\end{footnotesize}
\end{table}

\subsection{Exp-2b: Selector timing and accounting audit}
\label{app:defer-seltime}

The updated timing comparison toggles source deferral, not source pricing.
With deferral off, actions requiring both source and destination anchors
use maximum-weight closure solved by min-cut; single-anchor rounds can
still use DP. With deferral on, every action has one immediate anchor.
Table~\ref{tab:selector-real} covers eight full native graph streams.
The hardware matches the setup; solver and driver timing boundaries are
distinguished below.

\begin{table}[tb]\centering
\caption{\textbf{Exp-2b.} Updated real-stream selector timings.
Medians over three detector seeds at $\lambda=10^{-1}$.
Selection columns are full-stream totals; driver columns are per-batch
latencies including index emission. All times are in milliseconds.}
\label{tab:selector-real}
\label{tab:defer-seltime}
\begin{footnotesize}
\setlength{\tabcolsep}{0.3em}
\begin{tabular*}{\linewidth}{@{\extracolsep{\fill}}lrrrr@{}}
\toprule
& \multicolumn{2}{c}{Selection total} & \multicolumn{2}{c}{Driver / batch} \\
\cmidrule(lr){2-3} \cmidrule(lr){4-5}
Graph & off & on & off & on \\
\midrule
Facebook & $0.00$ & $0.00$ & $0.46$ & $0.46$ \\
Cit-HepTh & $0.05$ & $0.03$ & $2.34$ & $2.25$ \\
Cit-HepPh & $0.02$ & $0.01$ & $1.94$ & $2.18$ \\
Epinions & $2.07$ & $0.19$ & $4.66$ & $4.85$ \\
Twitter & $0.01$ & $0.00$ & $2.99$ & $3.01$ \\
DBLP & $0.71$ & $0.15$ & $6.23$ & $6.15$ \\
Flickr & $19.0$ & $1.90$ & $9.69$ & $9.62$ \\
YouTube & $155.4$ & $3.99$ & $27.2$ & $26.9$ \\
\bottomrule
\end{tabular*}
\end{footnotesize}
\end{table}

All $2{,}079$ reported enabled-arm rounds use DP. The supplied pooled
min-cut timing is $0.6455$\,ms over $804$ disabled-arm rounds; the reported
pooled DP timing is $0.0085$\,ms. Their approximately $76\times$ ratio is
descriptive, not a matched-input solver speed-up: the pools comprise
different optimization instances. The same-instance evidence is instead
Figure~\ref{fig:knee}(c), whose timings Table~\ref{tab:selector-exact} lists.
Facebook and Twitter contribute no informative
nonempty selection workload according to the supplied report; near-zero
rounded timings are not used to claim percentage improvements.

% Moved from the main text (2026-10-03); plotted as Figure fig:knee(c).
\begin{table}[tb]\centering
\caption{\textbf{Exp-2b.} Frozen-tree selector runtime.}
\label{tab:selector-exact}
\begin{footnotesize}
\setlength{\tabcolsep}{0.4em}
\begin{tabular*}{\linewidth}{@{\extracolsep{\fill}}crrrrr@{}}
\toprule
& \multicolumn{3}{c}{Time per solve (ms)} & \multicolumn{2}{c}{Speed-up over DP} \\
\cmidrule(lr){2-4} \cmidrule(lr){5-6}
$n$ & Enumeration & Closure & DP & Enumeration & Closure \\
\midrule
$\le 8$  & $0.0021$ & $0.0007$ & $\mathbf{0.0004}$ & $5\times$ & $1.7\times$ \\
$\le 12$ & $0.0242$ & $0.0010$ & $\mathbf{0.0004}$ & $54\times$ & $2.3\times$ \\
$\le 16$ & $0.1927$ & $0.0012$ & $\mathbf{0.0005}$ & $364\times$ & $2.3\times$ \\
$\le 20$ & $3.9625$ & $0.0018$ & $\mathbf{0.0009}$ & $4{,}615\times$ & $2.1\times$ \\
$\le 22$ & $13.681$ & $0.0021$ & $\mathbf{0.0011}$ & $12{,}699\times$ & $2.0\times$ \\
\midrule
$50$      & --- & $0.0076$ & $\mathbf{0.0027}$ & --- & $2.8\times$ \\
$200$     & --- & $0.0316$ & $\mathbf{0.0108}$ & --- & $2.9\times$ \\
$1{,}000$ & --- & $0.1816$ & $\mathbf{0.0500}$ & --- & $3.6\times$ \\
$5{,}000$ & --- & $1.0172$ & $\mathbf{0.2873}$ & --- & $3.5\times$ \\
$20{,}000$& --- & $5.6968$ & $\mathbf{1.2407}$ & --- & $4.6\times$ \\
\bottomrule
\end{tabular*}
\end{footnotesize}
\end{table}

\stitle{Timing boundary.} Selection counters accumulate solver-path time
over the entire stream, while driver latency is reported per batch and
includes graph maintenance, certificate upkeep, and index emission. The
driver denominator differs from Table~\ref{tab:main}'s maintenance-only
latency, which excludes partition output. Selection totals cannot be
divided directly by per-batch driver times to obtain a phase share.
The driver point estimates in Table~\ref{tab:selector-real} show no
comparable overall speed-up; no significance or equivalence claim is made.

\stitle{Corrected cost boundary.} The separate earlier paired cost check covers eight graphs
and all three detector instances over their full native streams. It uses
the driver's emitted refresh plan (\texttt{resum.tsv}) and the final pending
balance, i.e., $\mathsf K_T$ in Eq.~\ref{eq:cumulative-ledger}, including
unsettled work rather than assigning it zero cost. The two arms have identical
cost on $23$ of the $24$ graph--instance pairs; Epinions/\Louvain{} differs by
$+0.01\%$, and every pair agrees in $Q$ to five decimal places.

The earlier scorer (\texttt{hier\_cost.py}) reconstructed refresh work from
partition changes. It charged a removal-only source whenever its member set
changed, even if the driver had deferred its regeneration, producing apparent
increases of up to $44.8\%$. The paired deferral comparison therefore uses
schedule-plus-balance accounting. Table~\ref{tab:ablation} instead retains
the archived partition-derived structural scores and single-run timings;
its cost column is not a recomputed execution-plus-balance ledger.
The archived ablation runner launches six single-threaded jobs concurrently,
so these timings are descriptive rather than isolated-run measurements.
% Temporarily hidden with Exp-2c (2026-09-23); restore with tabepssweep.tex.
% Table~\ref{tab:eps} instead reports executed schedule cost and retained
% membership entries separately; its cost excludes the closing pending-share balance.
Comparing the component ablation's
internal control with Table~\ref{tab:main} is not a controlled defer-on/off
experiment: the main table already enables deferral.
The paired comparison above is the evidence for deferral's cost effect.

\subsection{Exp-3: Batch size and ingestion rate}
\label{app:exp3-detail}

We sweep batch sizes $E\in\{1,10,100,1000\}$ on MultiHop-RAG
($24.2$K edges; stream/base ratio $.91$) and Cit-HepPh
($420.9$K edges; stream/base ratio $.11$). Frontier labels in
Figure~\ref{fig:exp3} identify $E$.
The measured extraction rate is $r_{\mathrm{in}}=6.04$ edges/s. Under a fixed
batch size $E$, the plotted continuous-arrival approximation is
$E/(2r_{\mathrm{in}})$; summary generation proceeds at $0.945$ reports/s.
For fixed equally spaced arrivals, the exact mean batch-closing wait is
$(E-1)/(2r_{\mathrm{in}})$, including zero wait for $E=1$.
The plotted waiting time excludes subsequent maintenance and report generation,
and is not an end-to-end freshness measurement.

With $\lambda=0$, increasing $E$ from $1$ to $1000$ lowers MultiHop-RAG's
structural cost per edge from $113.5$ to $7.74$, while mean batch-closing delay
rises from $0.08$ to $82.8$\,s. Larger batches therefore amortize work but delay
index freshness. This trade-off is relevant to
delay-tolerant document-management workloads, where $E$ can be chosen from an
explicit latency budget. With the default $\lambda=10^{-1}$, however, cost-aware
selection removes most uneconomical openings before batching can amortize them:
singleton processing already costs $0.611$ per edge. This reduces the
incentive to use large batches solely to amortize refresh work.

Lazy evaluation further favors small batches when updates are localized. It
reduces singleton cost by $43\%$ on Cit-HepPh, whereas its effect is negligible
when the MultiHop-RAG increment perturbs most of the graph. Thus, lazy evaluation
complements cost-aware selection precisely when the update stream leaves enough
state reusable.

\subsection{Exp-4: Exchange-rate sweep}
\label{app:exp4-detail}

Figure~\ref{fig:knee} sweeps
$\lambda \in \{0, 10^{-3}, 10^{-2}, 10^{-1}, 0.3, 1, 3, 10, \infty\}$ on the six
graphs with an archived sweep. The source records $20$ batch iterations
per graph, a different horizon from the native stream lengths in
Table~\ref{tab:datasets}. We report relative structural cost
$\mathrm{Cost}(\lambda)/\mathrm{Cost}(0)$ and modularity loss
$Q(0)-Q(\lambda)$ over that recorded horizon;
Proposition~\ref{prop:mono} concerns a fixed frozen round and does not imply
monotonicity across the resulting trajectories.

All six curves are nearly flat from $\lambda=10^{-2}$ on; Cit-HepPh, DBLP,
and Flickr are already close to their floor at $\lambda=10^{-3}$.
Increasing the rate by a further three orders of magnitude, to $\lambda=10$,
changes normalized structural cost by at most $0.0045$ and modularity by less than
$10^{-4}$. At the knee, the cost term removes about $86$--$89\%$ of the
structural re-summarization work on five graphs
($\mathrm{Cost}(10^{-2})/\mathrm{Cost}(0)=0.110$--$0.137$), with modularity
losses of $0.06$--$0.67\times10^{-3}$. Thus, the location of the knee is stable
even though its magnitude varies by graph.

Facebook is the main cost-side exception: its normalized structural cost falls to $0.402$.
It has only $4{,}039$ vertices but average degree $43.7$, producing relatively
few, large communities and leaving less avoidable re-summarization work.
Epinions is the quality-side exception: it loses $0.67\times10^{-3}$ modularity,
compared with at most $0.19\times10^{-3}$ elsewhere, and still has normalized
structural cost $0.804$ at $\lambda=10^{-3}$.
The three reported Epinions values refer to different configurations:
$0.441$ is the full-graph
from-scratch \Louvain{} descriptor, $0.445$ is the from-scratch/reuse-aware
Leiden reference, and $0.458$ is the cost-oblivious sweep point here.

We use $\lambda=10^{-1}$, one decade beyond the common knee, as a shared default
rather than tuning it per dataset. A deployment may instead choose $\lambda$
from held-out update batches; the observed plateau confines the useful search
to the neighborhood of $10^{-2}$--$10^{-1}$.

\subsection{Post-hoc fidelity validation of structural weights}
\label{app:proxy-cal}

% Moved from Section 7.4 (2026-10-05).
\begin{table}[tb]
\centering
\caption{Structural weights versus measured summarizer input tokens.}
\label{tab:proxy-validation}
\begin{footnotesize}
\begin{tabular*}{\linewidth}{@{\extracolsep{\fill}}lrr@{}}
\toprule
\textbf{Diagnostic} & \textbf{Retained} & \textbf{Rebuilt} \\
\midrule
Leaf / internal samples & $944/91$ & $1{,}002/162$ \\
Leaf $R^2$ & $.851$ & $.859$ \\
Internal $R^2$ & $.380$ & $.437$ \\
\midrule
All-node $R^2$, $r_{\mathrm{int}}=1$ & $.222$ & $.143$ \\
All-node $R^2$, rescaled & $.675$ & $.738$ \\
Spearman $\rho$, $r_{\mathrm{int}}=1$ & $.746$ & $.623$ \\
Spearman $\rho$, rescaled & $.913$ & $.931$ \\
Fitted $r_{\mathrm{int}}$ & $3.41$ & $4.43$ \\
\bottomrule
\end{tabular*}
\end{footnotesize}
\end{table}

The benchmark's structural re-summarization cost is defined by entity and
child counts, not by a fitted token predictor. We validate its relationship to
measured summarizer inputs on the retained and fresh-rebuild versions of the
$327$-document MultiHop-RAG index. For each call in
\texttt{cache/community\_reporting}, we read
\texttt{usage.prompt\_tokens}. Because the cache key is a content hash, we map a
generated report back to \texttt{community\_reports.parquet} through its title
and retain only uniquely titled communities. This recovers $1{,}035/1{,}300$
communities ($80\%$) from the retained index and $1{,}164/1{,}300$ ($90\%$) from
the fresh-rebuild index; the diagnostic is limited to these matched samples.

\stitle{Fidelity and its limits.} We fit $\mathrm{tokens}=a+bn$ separately
to leaf entity counts and internal child counts. Table~\ref{tab:proxy-validation}
shows strong leaf fits and substantially weaker internal fits. Leaves dominate
the matched node counts, but these proportions do not measure token or refresh
cost mass. The leaf slopes are $78.4/81.8$ tokens per unit and the internal
slopes $267.3/362.9$. Their ratio gives the diagnostic multiplier
$r_{\mathrm{int}}=b_{\mathrm{internal}}/b_{\mathrm{leaf}}$; it is distinct from
the ledger ownership fraction $\alpha_{A,c}$. Post-hoc rescaling improves
all-node fit and rank agreement, but neither the remaining prediction error nor
the index-dependent multiplier supports a universal token conversion.
The rescaled correlations are validation results, not the fidelity claimed
for Table~\ref{tab:main}'s unchanged $r_{\mathrm{int}}=1$ weights.

\stitle{Relevance to selection.} The selector trades modularity gain against
structural cost through $g-\lambda w$, rather than predicting an absolute token
bill. A common multiplicative cost scale can be absorbed into $\lambda$.
Rank agreement therefore supports the proxy's ability to distinguish cheaper
from more expensive nodes, but is not sufficient to guarantee identical
decisions: $g-\lambda w$ also depends on cost magnitudes. The exactness result
concerns the stated structural objective, not a token-optimal objective.
Exp-5 supplies separate measured summary-token and maintenance-time ledgers;
this validation does not establish an absolute-token budget guarantee.

\stitle{Sensitivity to internal weighting.} Let $L,I$ denote the leaf and
internal contributions of DyLo's fixed refresh trace, including closing
liabilities, and $L_T,I_T$ the full-index totals. Post-hoc normalized cost is
\[
  C_{\mathrm{norm}}(r)=\frac{L+rI}{L_T+rI_T}.
\]
The multiplier enters both numerator and denominator, which can attenuate its
effect but does not make it cancel. On Epinions, changing $r_{\mathrm{int}}$
from $1$ to $3.41$ or $4.43$ changes DyLo's cost from $.132951$ to
$.137344$ or $.138955$, only $+3.3\%$ or $+4.5\%$. This one-sided sensitivity
check reweights only DyLo's fixed trace: baseline values, committed decisions,
$Q$, and latency are unchanged, and DyLo remains below every evaluated
incremental baseline in structural cost on this graph. It neither replaces
the main-table measurements nor asserts that every method ranking or selection
would survive re-optimization with new weights. We retain $r_{\mathrm{int}}=1$
as the common structural metric; Exp-5 separately measures native update-summary tokens and elapsed maintenance time.

\subsection{Exp-5: Matched-graph maintenance and sensemaking protocol}
\label{app:exp5-detail}

\stitle{Frozen workload and reuse.} The measured replay uses one archived
DyLo graph: $327$ documents, $13{,}264$ entities, $12{,}925$ relationships,
and $833$ text units, with the same $304$-document initial state and
$23$-document update for every arm. The final graph's semantic SHA-256 is
\texttt{cee9752bf403ba839}\allowbreak\texttt{0692cb8928ad5ce0}\allowbreak
\texttt{d3037983b3778a86}\allowbreak\texttt{fbf2caca8c196f9}.
DyLo replays the original \texttt{std4} binary and its archived partition,
using $\lambda=0.1$, $\sigma=0.2$, and hierarchy parameter $\gamma=5$.
The initial community reports are those of the shared archived DyLo index.
Baselines use their existing native adapters, including native partition
initialization where required. This is an update from one common archived
index, rather than a comparison of separately maintained per-method histories.
All incremental arms apply the archived lookup rule: take the first old
community with the same entity-membership set and reuse its report if present.
An input-content change without a membership
change does not trigger an additional refresh in this replay. We retain the
original materialization rule that filters communities without internal edges;
we do not add a singleton-coverage repair to this experiment. Complete report
coverage means every materialized community has a report row, not that every
extracted entity belongs to a retained root community. Rebuild regenerates
every materialized report. These choices reproduce the archived refresh rule
and may retain stale evidence in reused reports; they do not implement a
stricter evidence-change policy.

\stitle{Generation and timing.} Qwen2.5-7B-Instruct runs in BF16 with an
English structured report prompt requiring five findings, an output cap of
$2{,}048$ tokens, seed $17$, and eight concurrent requests. Response and
prefix caching are disabled for measured maintenance. Report source contexts
use a shared $4{,}000$-token \texttt{cl100k\_base} budget, while reported usage
uses Qwen's native tokenizer. The original three arms
use H100 GPU0 and the remaining five use H100 GPU1, with the same server
settings; resource-isolation logs are checked for every measured arm.
Embeddings use \texttt{BAAI/bge-small-en-v1.5}~\cite{xiao2024cpack}.
The graph-ready wall clock includes input loading, actual native detector
bootstrap and update, refresh selection, context preparation, summary
generation and retries, embedding, writing, and queryability validation.
It excludes initial index construction, model loading, raw-document extraction,
old-report repair, and all evaluation calls. We report one measured run per
arm, not repeated-run means or uncertainty estimates for timing.
\begin{table*}[tb]\centering
\caption{Graph-ready maintenance time by stage (percent of each arm's total).}
\label{tab:maintenance-phases}
\footnotesize
\setlength{\tabcolsep}{0.25em}
\begin{tabular*}{\linewidth}{@{\extracolsep{\fill}}lrrrrrr@{}}
\toprule
Method & Graph & Context & $\Delta$ reports & Merged reports & Publish & Other \\
\midrule
Rebuild & 0.04 & 0.91 & 0.00 & 98.66 & 0.34 & 0.04 \\
HIT-Leiden & 0.13 & 1.96 & 3.41 & 93.85 & 0.55 & 0.10 \\
\DynaMo{} & 0.15 & 1.92 & 3.32 & 94.01 & 0.51 & 0.10 \\
\QCA{} & 0.14 & 2.02 & 3.55 & 93.70 & 0.50 & 0.10 \\
DF-\Louvain{} & 0.12 & 1.93 & 3.40 & 93.97 & 0.49 & 0.10 \\
$\Delta$-screen & 0.11 & 1.91 & 3.29 & 94.06 & 0.52 & 0.10 \\
\Batch{} & 0.15 & 2.00 & 3.55 & 93.67 & 0.52 & 0.10 \\
\DyLo{} & 0.18 & 3.70 & 6.25 & 88.96 & 0.75 & 0.16 \\
\bottomrule
\end{tabular*}
\par\smallskip\raggedright
Graph includes native bootstrap and community maintenance. Publish includes embedding, index writing, and queryability checks. Other groups input loading, refresh selection, and residual overhead. Delta plus merged reports is the summary share in Table~\ref{tab:arena-main}.
\end{table*}

\stitle{Token accounting and exclusions.} Successful native prompt and
completion usage is added only for update-caused delta and merged community
reports. The $26$ delta reports are common to the incremental arms; Rebuild
directly summarizes the final index. Reports absent before the update are
repaired in a separately timed stage. The primary DyLo count is therefore
$394-50=344$, rather than a claim that only $344$ requests occurred in total.
\begin{table*}[tb]\centering
\caption{Separate old-report repair and failed update-summary calls.}
\label{tab:maintenance-exclusions}
\footnotesize
\setlength{\tabcolsep}{0.25em}
\begin{tabular*}{\linewidth}{@{\extracolsep{\fill}}lrrrrr@{}}
\toprule
Method & Repair calls & Repair (s) & Repair tok. (k) & Failed calls & Failed tok. (k) \\
\midrule
Rebuild & 0 & 0.0 & 0.0 & 0 & 0.000 \\
HIT-Leiden & 44 & 84.3 & 137.3 & 1 & 6.439 \\
\DynaMo{} & 39 & 74.7 & 116.9 & 0 & 0.000 \\
\QCA{} & 45 & 83.4 & 135.2 & 0 & 0.000 \\
DF-\Louvain{} & 37 & 77.4 & 107.3 & 1 & 6.431 \\
$\Delta$-screen & 37 & 67.9 & 102.5 & 0 & 0.000 \\
\Batch{} & 40 & 82.7 & 120.9 & 0 & 0.000 \\
\DyLo{} & 50 & 85.7 & 150.0 & 0 & 0.000 \\
\bottomrule
\end{tabular*}
\par\smallskip\raggedright
Repair tokens are successful native input plus output. Failed tokens are excluded from the primary successful-token total; retry time is retained in maintenance time.
\end{table*}
\begin{table*}[tb]\centering
\caption{Native root-report exposure, answer length, and successful maintenance token components.}
\label{tab:case-study-exposure}
\footnotesize
\setlength{\tabcolsep}{0.25em}
\begin{tabular*}{\linewidth}{@{\extracolsep{\fill}}lrrrrr@{}}
\toprule
Method & Root reports & Covered entities & Mean words & Input tok. (k) & Output tok. (k) \\
\midrule
Rebuild & 67 & 7,840 & 196.3 & 1926.118 & 1039.829 \\
HIT-Leiden & 67 & 7,841 & 205.7 & 1132.160 & 469.395 \\
\DynaMo{} & 71 & 7,841 & 207.4 & 1146.060 & 470.043 \\
\QCA{} & 67 & 7,840 & 198.8 & 1103.605 & 454.694 \\
DF-\Louvain{} & 69 & 7,841 & 218.3 & 1159.725 & 469.162 \\
$\Delta$-screen & 78 & 7,841 & 210.6 & 1174.995 & 478.749 \\
\Batch{} & 70 & 7,841 & 208.6 & 1125.332 & 463.547 \\
\DyLo{} & 151 & 7,521 & 214.0 & 918.156 & 239.070 \\
\bottomrule
\end{tabular*}
\par\smallskip\raggedright
Root cuts cover their materialized root entities, not every extracted entity. Input and output tokens here are community-maintenance usage, not answer-generation usage.
\end{table*}

There were $2$ failed update-summary calls consuming
$12,870$ native tokens in total, excluded from the successful ledger.
These failures reached the output cap and succeeded on retry.
Earlier format-repair probes and downstream question-answering calls are
outside maintenance accounting. The final answer-generation run records one
transport failure without returned usage; its token cost is unknown, and its
successful retry is retained. Preempted in-flight evaluation requests also
have unknown usage and are not assigned zero cost.

\stitle{Sensemaking, not reference-answer QA.} We use all $60$ archived
questions $q_0$--$q_{59}$, including questions receiving no-evidence answers.
They were frozen before generation on the new indices. The $15$-question
pilot's validated answers are retained for the first three arms, rather than
regenerated selectively. Each index supplies every report in its native
level-$0$ root cut to GraphRAG's global-search map--reduce implementation.
The same $8{,}000$-Qwen-token map and reduce evidence budgets apply to all
methods; map output is capped at $2{,}048$ tokens and reduce output at
$3{,}072$. Report counts and total query-token exposure are not matched across
partitions. Map points carry usefulness scores; the reducer synthesizes useful
points as English paragraphs. Structured outputs require a normal stop,
complete valid JSON, and allowed report identifiers; tool markers and truncated
outputs are rejected rather than stripped. Citations are rendered mechanically
without rewriting prose. All $480$ final answers are accepted only after
validation. Quality-generation time and tokens are outside maintenance totals.

\stitle{Independent judge and uncertainty.} Local \texttt{gpt-oss-120b}
uses temperature $0$, seed $17$, low reasoning effort, and a $2{,}048$-token
output cap, separately from the Qwen answer generator. We use the four
GraphRAG criteria~\cite{edge2024graphrag} and a blind pairwise prompt, permitting
ties. Each DyLo--baseline pair is assessed in both answer orders for each
question and criterion: $7\times60\times4\times2=3{,}360$ oriented cases.
Identical prompts reuse validated verdicts, so this is not a count of distinct
API calls. DyLo's win/tie/loss scores are $100/50/0$. We average both orders
within each question and then average the $60$ questions, using question-level
bootstrap intervals to preserve dependence between orders.
\begin{table*}[tb]\centering
\caption{DyLo preference scores and unadjusted $95\%$ question-bootstrap intervals.}
\label{tab:sensemaking-ci}
\footnotesize
\setlength{\tabcolsep}{0.25em}
\begin{tabular*}{\linewidth}{@{\extracolsep{\fill}}lrrrr@{}}
\toprule
Opponent & Comprehensiveness & Diversity & Empowerment & Directness \\
\midrule
\Batch{} & 60.00 [50.83, 69.58] & 66.25 [57.08, 75.42] & 53.33 [43.32, 63.75] & 49.58 [39.58, 59.58] \\
Rebuild & 56.67 [47.50, 65.42] & 61.25 [52.08, 70.42] & 48.33 [39.17, 57.50] & 48.33 [38.75, 57.92] \\
HIT-Leiden & 57.50 [47.92, 67.08] & 62.08 [51.25, 72.50] & 51.67 [41.25, 61.67] & 57.50 [47.08, 67.92] \\
\DynaMo{} & 64.17 [55.00, 72.92] & 63.75 [54.17, 73.33] & 52.50 [42.92, 61.67] & 55.00 [44.58, 65.42] \\
\QCA{} & 69.58 [60.83, 77.92] & 75.83 [66.25, 84.58] & 53.75 [43.33, 64.17] & 53.75 [43.75, 63.75] \\
DF-\Louvain{} & 55.00 [45.83, 63.75] & 56.25 [45.82, 66.67] & 52.50 [42.92, 62.08] & 47.08 [36.25, 57.92] \\
$\Delta$-screen & 53.33 [44.17, 62.92] & 60.83 [50.83, 70.83] & 48.75 [39.17, 58.34] & 40.00 [31.67, 48.75] \\
\bottomrule
\end{tabular*}
\par\smallskip\raggedright
$10{,}000$ resamples of the $60$ questions, seed $17$; both orders stay together within a question. Intervals are per comparison and criterion, without multiplicity correction.
\end{table*}

The predeclared process calibration passes: all four identical-answer controls
tie, no irrelevant padding is favored in eight controls, and directness
prefers the unpadded original in both orders. This small calibration does not
establish agreement with human judges or validate every refusal-versus-answer
case. Across the final comparisons, reversing answer order changes the
oriented preference on $13$--$29$ of $60$ questions per criterion.
Both orders are retained, rather than discarding inconsistent or unfavorable
cases. The reported intervals are unadjusted and do not prove a prespecified
non-inferiority margin. The $60$ questions are one fixed sample and answer
generation is single-run; this evaluation measures local-judge preferences,
not reference-answer accuracy or factual-grounding guarantees.

\stitle{Raw-document end-to-end scope.} The matched-graph replay does not
include fresh entity/relation extraction. A prior serial raw-document trial
used a different extracted graph and refresh policy and omitted two failed
text chunks; its timings cannot be added to this replay's maintenance results
to create a measured full-update comparison. The full-update expression in
Section~\ref{sec:exp:e2e} is conditional on a shared upstream ingestion stage
and its maintenance fraction. A matched raw-document-to-queryable comparison
is required to turn that implication into an observed end-to-end saving.

\subsection{Archived Exp-5: Logical input work and ranked judging}
\label{app:exp5-archived}

This archived experiment is retained for provenance. It uses independently
extracted graphs, incomplete report coverage, a different summary prompt and
output cap, the latter $60$ synthesis questions, and a different judge and
query-exposure protocol. Its logical tokens and mean ranks are not the native
usage or preference scores of the new matched-graph replay.
\begin{table}[tb]\centering
\caption{\textbf{Exp-5.} Archived logical input work and answer ranks with matched report counts.}
\label{tab:arena-archived}
\footnotesize
\setlength{\tabcolsep}{0.20em}
\begin{tabular*}{\linewidth}{@{\extracolsep{\fill}}lrrrrrr@{}}
\toprule
& \multicolumn{2}{c}{Input work} & \multicolumn{4}{c}{Mean rank ($\downarrow$)} \\
\cmidrule(lr){2-3} \cmidrule(lr){4-7}
Method & Targets & Tok. (k) & Comp. & Div. & Emp. & Dir. \\
\midrule
Rebuild         & $1{,}518$ & $5{,}493$ & $4.75$ & $4.75$ & $4.58$ & $\mathbf{3.81}$ \\
HIT-Leiden      & $985$ & $4{,}944$ & $5.44$ & $5.38$ & $5.66$ & $4.40$ \\
\DynaMo{}       & $825$ & $4{,}230$ & $4.64$ & $4.81$ & $4.53$ & $4.23$ \\
\QCA{}          & $829$ & $4{,}224$ & $4.31$ & $4.32$ & $4.44$ & $4.79$ \\
DF-\Louvain{}   & $821$ & $4{,}254$ & $3.92$ & $3.75$ & $3.53$ & $4.82$ \\
$\Delta$-screen & $818$ & $4{,}328$ & $5.17$ & $5.24$ & $5.23$ & $4.48$ \\
\Batch{}        & $777$ & $3{,}616$ & $4.72$ & $4.68$ & $4.60$ & $3.93$ \\
\midrule
\DyLo{} & $\mathbf{394}$ & $\mathbf{2{,}514}$ & $\mathbf{3.04}$ & $\mathbf{3.08}$ & $\mathbf{3.43}$ & $5.53$ \\
\bottomrule
\end{tabular*}
\par\smallskip\raggedright
Tok. (k): logical input tokens in thousands. Ranks use $472$ blind
rankings across four criteria ($1$ = best of $8$); details in
Appendix~\ref{app:exp5-archived}.
\end{table}


\stitle{Models and decoding.} Community summaries are generated with
\texttt{Qwen/Qwen2.5-7B-Instruct}~\cite{yang2024qwen25} using the report
model's \texttt{max\_tokens=8192}, a $4{,}000$-token source-context budget,
and a $2{,}000$-word report limit. The extraction model separately uses
\texttt{max\_tokens=2048}. Logical input-work accounting uses
\texttt{cl100k\_base} before cache lookup; it is not native billed usage. Embeddings use
\texttt{BAAI/bge-small-en-v1.5}~\cite{xiao2024cpack}, with inputs truncated to $512$ tokens. The
blind judge is Opus~5.

\stitle{Matched report exposure.} Each incremental partition seeds GraphRAG's
\texttt{hierarchical\_leiden}. All arms use a common \texttt{max\_cluster\_size}.
For answer generation, every arm's level-$\le1$ report pool is randomly
downsampled to $453$ under a fixed seed, and all arms use the same response
template with a $250$-word limit. This matches the number of available reports,
not their content or token volume. The logical input ledger counts source-context and template tokens before
cache lookup and before report-pool sampling. It excludes output tokens and
API-level retries. Final report coverage and extracted graphs are not matched
in these archived arms, unlike the new matched-graph replay.
The final DyLo, Batch, and Rebuild indexes lack $82$, $107$, and $136$
reports, respectively, so target counts are not successful-generation counts.

\stitle{Blind ranked judging.} For each of the $60$ questions, the judge ranks
the eight anonymized answers from best to worst, without ties, on each of the
four criteria~\cite{edge2024graphrag}. The letter-to-arm mapping is redrawn
independently for every question, which limits the position bias documented
for LLM judges~\cite{zheng2023judging}. The full set is graded twice under two
independent permutations, and only questions judged in both rounds are kept;
each round is split into two batches, and the four batches are judged in
mutually independent contexts. One question,
on which all eight arms returned the same refusal, carries no signal and is
excluded, leaving $59\times4\times2=472$ rankings. Mean ranks and pairwise win
rates are computed from this same pool. Pooled over comprehensiveness,
diversity, and empowerment, \DyLo{} wins $68.8\%$ of all pairwise comparisons
and $60$--$76\%$ against each baseline; DF-\Louvain{} is second at $61.0\%$,
and \QCA{}, \DynaMo{}, \Batch{}, and Rebuild lie within $46$--$54\%$ of one
another. HIT-Leiden and $\Delta$-screen decline $11.9\%$ and $8.5\%$ of
questions, versus $3.4$--$6.8\%$ for the other arms.

\stitle{Protocol checks.} Across the two permutations, pairwise orderings agree on
$90.5\%$ of arm pairs ($88.1$--$91.8\%$ per criterion). Comprehensiveness and
diversity agree on $93.5\%$ of pairwise orderings, so they should be read as
one finding rather than two; empowerment is more distinct ($82.4$--$83.7\%$)
and directness nearly independent of the other three ($30.7$--$44.0\%$).

DyLo's largest level-$\le1$ community contains $209$ entities, compared with
$282$--$707$ for the baselines, while method-level medians lie between $10$ and
$13$. The separate reference-answer experiment below
provides retrieval-read measurements, which are not measurements of this arena.

\subsection{Reference-answer quality at the final snapshot}
\label{app:groundtruth}

The reference-answer diagnostic uses the full $609$-article MultiHop-RAG
corpus in publication order, starting with $304$ articles and ingesting the
remaining $305$ in $20$ batches. The selected $414$ questions comprise $225$
with incremental evidence and $189$ controls. As described in
Section~\ref{sec:exp:e2e}, this is a separate community-summary retrieval
pipeline rather than the GraphRAG global-search arena, evaluated only at the
final snapshot. Qwen generates summaries from at most $6{,}000$ characters
of community evidence, with a $180$-word instruction and $350$-token output
limit. The $15$ nearest entities are retrieved with BGE embeddings;
communities are selected by entity-hit count, with at most six summaries and
$12{,}000$ characters supplied to Qwen for a $60$-token short answer.
Token $F_1$ uses normalized word overlap and the best score over references. Unchanged member sets inherit initial-state summaries;
changed communities use final-state evidence. This inheritance rule tests
the archived membership-based policy, not a certificate of unchanged evidence.
Table~\ref{tab:groundtruth-scores} reports the scores and paired
normal-approximation intervals, which are exploratory and not adjusted for
multiple comparisons. Zero-containing intervals do not establish
equivalence or non-inferiority.

\begin{table}[tb]\centering
\caption{Reference-answer check: token $F_1$ and paired DyLo gains on $414$ questions.}
\label{tab:groundtruth-scores}
\begin{footnotesize}
\setlength{\tabcolsep}{0.25em}
\begin{tabular*}{\linewidth}{@{\extracolsep{\fill}}lrr@{}}
\toprule
Method & $F_1$ & DyLo gain [95\% CI] \\
\midrule
\DyLo{}-Leiden & $.1866$ & --- \\
\QCA{} & $.1723$ & $.0143\;[-.0090,.0375]$ \\
FS \Louvain{} & $.1662$ & $.0204\;[-.0025,.0433]$ \\
FS Leiden & $.1570$ & $.0296\;[.0028,.0564]$ \\
\Batch{} & $.1546$ & $.0320\;[.0071,.0569]$ \\
\DynaMo{} & $.1397$ & $.0469\;[.0220,.0718]$ \\
\GreMod{} & $.0886$ & $.0980\;[.0683,.1277]$ \\
\bottomrule
\end{tabular*}
\par\smallskip\raggedright
Intervals are paired normal approximations
$\bar d\pm1.96s_d/\sqrt{414}$, unadjusted for multiple comparisons.
\end{footnotesize}
\end{table}

Report counts are similar: \DyLo{} reads $5.10$ summaries per query,
versus $5.14$ for \QCA{} and $4.45$--$5.14$ across all baselines.
These counts do not control retrieved-token volume or summary length.
% Temporarily hidden with the reference-answer control table, 2026-09-24.
% The same-partition controls in Table~\ref{tab:mhrag-controls} compare selective
% deferral with refreshing every summary and with blanking $40\%$ of summaries.

\section{Question set and corpus provenance}
\label{app:questions}

\stitle{Corpus.} The GraphRAG case studies use MultiHop-RAG~\cite{tang2024multihoprag},
the news collection GraphRAG itself evaluates on. It contains $609$ English news
articles published between 2023-09-26 and 2023-12-25 and drawn from $49$ outlets,
the largest being Sporting News ($101$ articles), TechCrunch ($97$), The Verge
($45$), Polygon ($44$), and The Guardian ($29$). By category, the corpus contains
sports ($211$), technology ($172$), entertainment ($114$), business ($81$),
science ($21$), and health ($10$) articles. For Exp-5, we index the first $304$ articles in
publication order and ingest the next $23$ as one batch, giving the $327$-document
arena index. Appendix~\ref{app:groundtruth} describes the separate full-corpus
reference-answer stream.

\stitle{Question generation.} MultiHop-RAG provides reference-answer queries,
which we use for the ground-truth freshness check in
Section~\ref{sec:exp:e2e}. The GraphRAG global-search evaluation~\cite{edge2024graphrag} does not
provide reference answers for open-ended corpus-level synthesis; its protocol
instead asks a language model to imagine potential corpus users, their tasks,
and the questions raised by each task, given only a prose corpus description.
We use the same corpus-description protocol to generate $120$ synthesis
questions. The new measured Exp-5 uses the first $60$, $q_0$--$q_{59}$;
the archived ranked experiment uses the latter $60$, $q_{60}$--$q_{119}$.
Both sets are listed below. Repeated judging does not increase the number of
questions. The generator never
sees an index, community partition, or system answer.

\subsection{Measured replay questions (\texorpdfstring{$q_0$--$q_{59}$}{q0-q59})}
\begin{footnotesize}
\begin{enumerate}[label={\(q_{\arabic*}\).},start=0]
\setlength{\itemsep}{1pt}
\item What are the most frequently mentioned entities across all categories from September to December 2023?
\item How has the frequency of articles on specific topics changed from September to December 2023?
\item Can you detect any patterns in the type of articles (e.g., opinion pieces vs. news reports) over the four months?
\item Which categories saw the highest increase in article count from September to December 2023?
\item What is the overall sentiment towards major corporations mentioned in the business category?
\item How does public sentiment towards a particular celebrity change from September to December 2023?
\item Can you determine if there is a correlation between positive mentions of a technology company and their stock performance?
\item What is the sentiment towards environmental issues as reported in the health and science category?
\item How do the most common themes in the entertainment category compare with those in the technology category?
\item What are the key differences in the types of entities discussed in sports versus business articles?
\item Can you identify any significant overlaps in the entities discussed in health and science articles compared to technology articles?
\item How does the frequency of articles about international events differ between the sports and business categories?
\item How consistent is the use of sources across different categories?
\item What is the distribution of original reporting versus rehashing of press releases or other media?
\item Can you assess the diversity of viewpoints presented in the articles?
\item How often do articles reference the same entities, and what does this imply about the scope of coverage?
\item How did major global events (e.g., political elections, natural disasters) influence the content of the articles across categories?
\item What are the top three most covered events or issues in each category, and how did these topics evolve over the four months?
\item Can you identify any long-term trends in the coverage of certain events that may indicate a shift in editorial focus?
\item How does the coverage of local events compare with that of global events in the sports and technology categories?
\item What themes have emerged consistently across different categories throughout the period?
\item How do the themes in the entertainment category compare with those in the health and science category?
\item Can you identify any thematic shifts in the coverage of business news from September to December 2023?
\item How do the thematic focuses of the sports category differ from those in the technology category?
\item What are the most popular topics among readers based on the frequency of mentions?
\item How do consumer interests vary between different categories?
\item What emerging trends can be identified from the data that could inform marketing strategies?
\item Can you determine which entities are most likely to drive consumer interest in the business category?
\item Which brands have seen the highest increase in visibility across the four months?
\item How does brand visibility differ across the various categories?
\item Can you identify any brands that have maintained consistent visibility throughout the period?
\item Which brands are most associated with positive sentiment in the technology category?
\item How does the coverage of a specific industry (e.g., healthcare, tech) affect its visibility and perception?
\item What are the most influential companies in terms of driving industry-wide discussions?
\item Can you determine if there is a correlation between increased coverage and changes in stock prices for listed companies?
\item How does the coverage of a particular industry compare with its actual market performance?
\item What are the most frequently discussed health-related issues during the period?
\item How do public health concerns in the news compare with those in medical journals?
\item Can you identify any trends in the coverage of new health technologies or treatments?
\item How does the coverage of mental health issues compare with that of physical health issues?
\item What is the public's sentiment towards specific health issues (e.g., obesity, diabetes)?
\item How does the coverage of a particular health issue affect public awareness and behavior?
\item Can you determine if there is a correlation between the volume of articles and public health outcomes?
\item How does the coverage of a specific health issue (e.g., vaccination campaigns) vary across different regions?
\item What are the most discussed technological advancements in healthcare during the period?
\item How do these technological advancements compare with their actual adoption rates in clinical practice?
\item Can you identify any trends in the integration of technology into healthcare practices?
\item How does the coverage of telemedicine compare with traditional healthcare methods?
\item What are the most frequently discussed technological advancements during the period?
\item How do these technological advancements compare with the latest innovations in the tech industry?
\item Can you identify any emerging trends in the tech sector that could present new opportunities?
\item How does the coverage of a specific technology (e.g., AI, blockchain) compare with its actual adoption rate?
\item What are the most discussed tech companies during the period, and how do they compare with each other?
\item How does the coverage of a specific tech company (e.g., Apple, Google) affect its market position?
\item Can you determine if there is a correlation between increased coverage and market share growth for specific tech companies?
\item How does the coverage of a particular tech trend (e.g., cybersecurity) vary across different industries?
\item What is the public's sentiment towards specific technologies (e.g., smartphones, smart homes)?
\item How does the coverage of a particular technology (e.g., electric vehicles) affect consumer perception and behavior?
\item Can you identify any trends in the adoption of new technologies based on the coverage?
\item How does the coverage of a specific technology (e.g., virtual reality) compare with its actual usage rates?
\end{enumerate}
\end{footnotesize}

\subsection{Archived synthesis questions
(\texorpdfstring{$q_{60}$--$q_{119}$}{q60-q119})}
\begin{footnotesize}
\begin{enumerate}[label={\(q_{\arabic*}\).},start=60]
\setlength{\itemsep}{1pt}
\item How has the frequency of coverage on specific topics changed from September to December?
\item What are the most common themes across different categories of news articles?
\item Which entities have been covered consistently across multiple categories, and why might this be significant?
\item Identify any emerging or declining trends in the types of stories being reported.
\item Can you identify any biases in the reporting based on the tone and language used?
\item How does the coverage of major events differ between categories?
\item Are there any discrepancies in the portrayal of certain entities across different news articles?
\item Examine if there is a correlation between the source of information and the angle of reporting.
\item How have global events influenced the local news coverage during the period?
\item Identify any news articles that were directly affected by international developments.
\item Compare the impact of domestic versus international news events on the local audience.
\item Determine if there are any patterns in how different categories of news articles respond to the same global event.
\item Calculate the overall sentiment score for each category of news articles over the period.
\item Identify any significant shifts in sentiment scores for specific entities over time.
\item Compare the sentiment scores of different entities within the same category to see if there are any notable differences.
\item Analyze the relationship between the volume of news articles and the overall sentiment score.
\item Using clustering algorithms, group similar articles together based on their content and topics.
\item Determine the number of clusters that best represent the diversity of the dataset.
\item Evaluate the coherence of the clusters formed, ensuring they align with expected categories.
\item Identify any outliers or anomalies within the clusters and explain their presence.
\item Based on the historical data, predict which topics are likely to gain more attention in the next quarter.
\item Assess the predictive power of certain features (e.g., entity mentions, article length) in forecasting future trends.
\item Evaluate the accuracy of your predictions using a validation set of recent articles.
\item Discuss any challenges in predicting future trends based on the current dataset.
\item Which demographic segments are most likely to engage with the news articles in each category?
\item Analyze the content to determine which topics are most appealing to different age groups.
\item Identify any gaps in the current coverage that could attract new audiences.
\item Determine the effectiveness of the current targeting strategies based on the engagement levels.
\item Which types of articles tend to perform better on social media platforms?
\item Analyze the engagement metrics to identify the most effective headlines and formats.
\item Determine the optimal posting times for maximizing reach and engagement.
\item Evaluate the impact of multimedia elements (e.g., images, videos) on article performance.
\item Based on the dataset, what content themes should be prioritized in future articles to maintain high engagement?
\item Identify any underrepresented topics that could be included to diversify the content portfolio.
\item Assess the balance between different categories to ensure a well-rounded content strategy.
\item Develop a plan for incorporating trending topics into the content calendar.
\item Track changes in public opinion on specific issues over the period based on the articles.
\item Identify any shifts in public sentiment towards certain entities or topics.
\item Analyze the role of influential figures in shaping public opinion through the articles.
\item Determine the effectiveness of news articles in influencing public opinion on key issues.
\item Examine the extent to which news articles have influenced policy decisions during the period.
\item Identify any instances where news coverage led to policy changes or public debates.
\item Assess the role of media in highlighting or downplaying certain issues in the policy agenda.
\item Discuss the implications of media influence on policy-making processes.
\item Analyze the relationship between the news coverage and societal trends during the period.
\item Identify any correlations between news articles and changes in societal attitudes or behaviors.
\item Explore the impact of news articles on public discourse and cultural norms.
\item Discuss the broader implications of the findings for understanding the role of media in society.
\item Review the consistency and quality of reporting across different categories.
\item Identify any areas where the reporting standards could be improved or standardized.
\item Assess the accuracy and reliability of the information presented in the articles.
\item Determine the effectiveness of fact-checking processes in maintaining editorial integrity.
\item Analyze the factors that contribute to higher reader engagement in different categories.
\item Identify strategies to increase the reach and impact of news articles.
\item Determine the most effective ways to integrate multimedia elements into the articles.
\item Evaluate the impact of interactive features (e.g., polls, quizzes) on reader engagement.
\item Assess the balance between different types of content (e.g., opinion pieces, investigative reports) in the current mix.
\item Identify any gaps in the current content mix that could be filled to better serve the audience.
\item Determine the most effective ways to incorporate user-generated content into the news articles.
\item Evaluate the potential for partnerships with other media outlets to enhance the content mix.
\end{enumerate}
\end{footnotesize}
